% !PRINTED 491
\weberpage{491}
%<<CONT
\gesperrt{; its coefficients $A_1, A_2, \ldots, A_m$ are integers.}

We remark that, according to this definition, in order to characterise an algebraic number $\Theta$ as integral, it suffices if among the infinitely many equations of the form (1) which $\Theta$ satisfies there is \gesperrt{one} whose coefficients are integers.

The integral algebraic numbers include as a special case the ordinary integers, which for distinction we call \gesperrt{integral rational numbers}. The positive integral rational numbers we also call, following a widespread usage, \gesperrt{natural numbers}.

By integers without qualification we then understand integral algebraic numbers, rational as well as irrational.

1. \gesperrt{An integral algebraic number which is at the same time rational is necessarily an integral rational number.}

For suppose that $\Theta = P:Q$ is a rational fraction, and that $P, Q$ are integral rational numbers without common divisor, say $Q$ positive; then from (1) there results
\[
P^m + A_1 P^{m-1} Q + A_2 P^{m-2} Q^2 + \cdots + A_m Q^m = 0,
\]
and from this it is to be seen that every prime divisor of $Q$ would have to be contained in $P$. Hence we must have $Q = 1$, and $\Theta = P$ is an integral rational number.

2. \gesperrt{The sum, difference and product of two integers are again integers.}

To prove this principal theorem, we assume that $\alpha, \beta$ are two integers which satisfy the equations
\weq{2}{\begin{array}{l}
\alpha^\mu + a_1 \alpha^{\mu-1} + \cdots + a_{\mu-1} \alpha + a_\mu = 0 \\
\beta^\nu + b_1 \beta^{\nu-1} + \cdots + b_{\nu-1} \beta + b_\nu = 0
\end{array}}
and make one of the three assumptions
\[
\omega = \alpha + \beta, \ \alpha - \beta, \ \alpha \beta.
\]
Then we set $\mu \nu = m$ and denote the $m$ quantities
\[
\alpha^r \beta^s \quad \begin{array}{l} r = 0, 1, \ldots, \mu - 1 \\ s = 0, 1, \ldots, \nu - 1 \end{array}
\]
in any order by $\omega_1, \omega_2, \ldots, \omega_m$.
